% University of Helsinki
% Department of Mathematics and Statistics
% Template for Bachelor's or Master's thesis
%
% This template is intended as a starting point for a thesis;
% feel free to customize it or use your own template.
% See the comments for hints and customization.
%
% Initial version: 29 August 2024 / Petri Laarne


% Add the 'twoside, openright' options when producing a proper printed version
\documentclass[a4paper, 11pt]{report}

\usepackage[hidelinks]{hyperref} % Always load this: clickable links and PDF reader navigation bar
\usepackage{lastpage} % Used to determine the total page count
\usepackage{calc} % For layout calculations involving TeX length units (template title page)
\usepackage{graphicx} % For embedding image files (JPEG, PNG, PDF)
\usepackage{subcaption} % For subfigures
\usepackage{tikz} % For drawing diagrams and other complex vector graphics

% If you use the LuaLaTeX or XeLaTeX compiler, comment out the following line.
% These compilers offer better support for Unicode, so you should use them if possible.
\usepackage[T1]{fontenc}
\usepackage[most]{tcolorbox}
%\DeclareAutoCiteCommand{textcite}{\textcite}{\textcites}
%\ExecuteBibliographyOptions{autocite=textcite}
\DeclareMathOperator*{\argmax}{arg\,max}

\newtcolorbox[auto counter]{proposition}[1][]{
  enhanced,
  breakable,
  fonttitle=\scshape,
  title={Proposition \thetcbcounter},
  #1
}

% Pass the language of the thesis as the LAST option in the list below.
% This changes the language used for dates and words like 'Chapter' or 'Bibliography'.
% You can remove the languages that do not appear in your work.
\usepackage[finnish,swedish,english]{babel}

\usepackage{csquotes} % This package is suggested by 'babel'

% Bibliography. Change the latter command if you modify the .bib file name.
\usepackage{biblatex}
\addbibresource{Resources/library.bib}

\usepackage{mathtools} % Extended mathematics support (this package extends amsmath)
\usepackage{amssymb} % More mathematical symbols
\usepackage{amsfonts} % Blackboard bold mathematics font

% Theorem environments
\usepackage{amsthm}
\theoremstyle{plain}
\newtheorem{theorem}{Theorem}[section] % Theorems are numbered within chapter and section
\newtheorem{lemma}[theorem]{Lemma} % Lemmas etc. follow the same numbering as theorems
\theoremstyle{definition}
\newtheorem{definition}[theorem]{Definition}
\theoremstyle{remark}
\newtheorem{remark}[theorem]{Remark}

% Some examples of custom mathematical commands
\newcommand{\ZZ}{\mathbb Z} % Blackboard bold Z
\newcommand{\RR}{\mathbb R} % Blackboard bold Z
\newcommand{\dx}{\mathop{}\!\mathrm d x} % Nice spacing (and upright d) for dx in integrals

% Finnish notation for evaluating integrals
% Instead of \int_a^b syntax you must use it as \sij{a}{b}
% (To enable the sub-/superscript syntax would require some more complicated TeX wizardry.)
\newcommand{\sij}[2]{\mathop{\bigg/}\nolimits_{\hspace{-0.65em}#1}^{\hspace{0.05em}#2}\hspace{-0.2em}}


\begin{document}

%
% TITLE PAGE
%
\begin{titlepage}

\noindent
\begin{minipage}[b]{\textwidth-45pt}
University of Helsinki\\
Faculty of Science\\
Bachelors Program in Mathematics
% Helsingin yliopisto\\
% Matemaattis-luonnontieteellinen tiedekunta\\
% Helsingfors universitet\\
% Matematisk-naturvetenskapliga fakulteten\\
\end{minipage}
\hfill\includegraphics[width=45pt]{HYlogo.pdf}

\vspace{4pt}\hrule\vfill

\begin{center}
Bachelor's thesis\\[8pt]
% Kandidaatintutkielma / Maisterintutkielma\\[8pt]
% Kandidatavhandlingen / Magisteravhandlingen\\[8pt]

{\huge\bfseries Title of the thesis}\\[8pt]

Santeri Väätäjä
\end{center}

\vfill\hrule\vspace{4pt}

\noindent
Supervisor:\hfill Month and year of submission\\
Petri Ola\\
\end{titlepage}

%
% ABSTRACT
%
% This command starts a new page.
% In a print version ('twoside' option passed to \documentclass), the page is on the right side
\pagenumbering{gobble}
\cleardoublepage

\noindent\textbf{Title:} Repeat the title here\\ % Otsikko / Titel
\textbf{Author:} Santeri Väätäjä\\ % Tekijä / Författare
\textbf{Month and year:} Month and year of submission\\ % Aika / Datum
\textbf{Page count:} \pageref*{LastPage}~pp.\\[1em] % Sivumäärä: XX~s. / Sidoantal: XX~s.
% The page count is updated automatically, but it might require two compilations to be correct.
% If you customize the backmatter, this might become incorrect - check and correct manually if needed.

% If you include an abstract in a different language, wrap it in a foreignlanguage block
\begin{foreignlanguage}{english}
\noindent\textbf{Abstract:}\\ % Tiivistelmä / Referat
The abstract should be about one page long.
It also doubles as the maturity test.
Do not use mathematical notation in the abstract,
since it will also be displayed on the electronic thesis repository.

For the publication date, put in the month and year of submission of the finished thesis.
Also include 3~to 5~keywords that support searching for your thesis.
They are also shown one the repository page.

You can include the abstract in more than one language by copying and modifying
this code between comment blocks.
To get correct hyphenation for the language (and not with the rules of the main language),
put the text in a \verb|foreignlanguage| block as shown in the code
(just modify the language parameter to match the abstract).
\end{foreignlanguage}


\vfill

\noindent\textbf{Keywords:} Fill in some keywords, Separated by commas % Avainsanat / Nyckelord




%
% TABLE OF CONTENTS
%
\cleardoublepage
\tableofcontents


%
% MAIN TEXT STARTS HERE
%

% The \pagenumbering{arabic} command must be immediately after the first chapter title.
% Do not repeat it for other chapters.
\chapter{Introduction}
\pagenumbering{arabic}

\chapter{Poisson Kernel in $\RR_+^{n+1}$}
In order to define the Poisson kernel

%\begin{definition}[Summability to $L$]\label{def:mean_of_int}
%    A function $f$ is said to be $M$-summable to $L$, if $\lim_{\epsilon\rightarrow0}M_{\epsilon,\Phi}(x)=L$, where
%
%    \begin{equation*}
%        M_{\epsilon,\Phi}(f)=M_{\epsilon}(f)=\int_{\RR^n}\Phi(\epsilon x)f(x)dx
%    \end{equation*}

%    \noindent where $\Phi\in C_0$ and $\Phi(0)=1$. Then $M_\epsilon(f)$ is called the $\Phi$-mean of the integral.
%\end{definition}

\clearpage
\chapter{Harmonic functions in $E_+^{n+1}$}

\begin{theorem}
    If $\mu\in M(E_+^{n+1})$ and $u(x,y)=\int_{E_n}P(x-t,y)d\mu(t)$ is the Poisson-Stiltjes integral of $\mu$ then $u$ is harmonic in $E_+^{n+1}$ and 

    \begin{equation}
        ||u(\cdot,y)||_1=\int_{E_n}|u(x,y)|dx\leq||\mu||
    \end{equation}

    \noindent where $||\mu||$ is the total variation, $|\mu|(E_n)$, of $\mu$. Furthermore,

    \begin{equation*}
        \lim_{y\rightarrow0}\int_{E_n}u(x,y)\varphi(x)dx=\int_{E_n}\varphi(x)d\mu(x)
    \end{equation*}

    \noindent for all $\varphi\in C_0$; that is, $u(\cdot,y)$ converges to $\mu$ in the weak* topology.
\end{theorem}

\clearpage
\chapter{Nontangential Convergence of Harmonic Functions}


%
% BIBLIOGRAPHY
%
% The [heading=bibintoc] option adds the bibliography to table of contents
\printbibliography[heading=bibintoc]




%
% APPENDICES (if necessary)
%
%\appendix

%\chapter{Some more information}

%If you need to have appendices in your thesis,
%you can write them as usual chapters after the \verb|\appendix| command.
%Appendices are numbered alphabetically.

\end{document}
